\begin{afterword}
\summaryhead

Written the day after ``Tsuyoku Naritai!'', the post opens with hunter-gatherers, who are usually egalitarian and whose successful hunters downplay their kills to avoid envy. Anyone who aims to do their best, the author argues, will sooner or later aim above the average, and a person ashamed of doing better than others will stop at the median.

The author admits finding pride in beating others a useful motivator, ``despite my principles,'' and compares such competition to Go. In any case, it is not healthy to be ashamed of doing better, and life is not graded on a curve: run as fast as you can without worrying about pulling ahead. Progress will leave you with fewer sins to confess. It may be wise to downplay your success in public, but you should keep whispering ``Tsuyoku, tsuyoku!'' to yourself, set a higher target, and keep running whether ahead or behind.

\responsehead

The advice is modest and sensible: do not be ashamed of improving, and if public modesty is wise, practice it without letting it limit what you aim at. The difficulty is that the post gives no evidence of the problem it sets out to solve.

The opening example is sound as anthropology. Hunter-gatherers do tend to be egalitarian, and Richard Lee's account of the !Kung describes the hunter who says he saw ``just a little tiny one.'' But in Lee's account the modesty is enforced by the others, to keep a good hunter from pride; the hunter's success is not in doubt, only his boast. What the example shows is public downplaying, which the post itself later calls possibly ``the course of wisdom.'' What the post argues against is something else: private shame at doing better, which makes ``the median'' a ``concrete wall.'' For that it offers no example and no evidence. The title calls it an instinct, and the only link from the foragers to the reader is the phrase ``the ancestral environment.''

The post's two defenses of doing better also do not fit together. Pride in beating others is ``a useful motivator''; and besides, ``life is not graded on a curve,'' so run ``without worrying that you might pull ahead.'' The first draws on comparison with others; the second says not to worry about it. The post is candid that the first goes against ``my principles,'' which are probably those of ``Twelve Virtues of Rationality,'' the source of the second. It does not say how the two fit together.

In short: a pep talk against being ashamed of doing better, opened with a hunter-gatherer custom of public modesty, which the post itself allows, rather than with any case of the private shame it argues against.
\end{afterword}
